% Section content template for individual Educative lessons
% This file contains the actual content for a specific section
% Generated from Educative API section content

% Section: Requirements for the Amazon Locker Service
% Section ID: 5801016767545344
% Section Slug: requirements-for-the-amazon-locker-service
% Generated: 2025-12-11T15:16:31.013372

This lesson outlines the functional and operational requirements for the Amazon Locker service. Identifying and understanding requirements is essential to defining the system's scope and ensuring a robust, user-friendly design.
\\
We'll use the notational convention to identify each requirement with a unique label ``Rn,'' where ``R'' is short for requirement and ``n'' is a natural number.

\subsection{Requirement collection}\label{slaPtV0hrXKeJxBPF9gA3}
\\
The requirements for the Amazon Locker service are defined below:

% \vspace{-1.0\baselineskip}
\noindent
\begin{minipage}[t]{0.48\textwidth}\vspace{0pt}
\begin{itemize}
\item
\textbf{R1:} Customers can select a preferred locker location for order pickup during checkout.
\item
\textbf{R2:} An order may contain one or more items. Based on locker size availability, items are packaged together if possible.
\item
\textbf{R3:} Locker locations contain multiple lockers of various sizes (extra small, small, medium, large, extra large, double extra large).
\item
\textbf{R4:} Only packages that fit fully within the locker's interior dimensions are eligible for delivery.
\end{itemize}
\end{minipage}
\hfill
\begin{minipage}[t]{0.48\textwidth}\vspace{0pt}
\textit{Column component type 'MxGraphWidget' not yet supported.}
\end{minipage}

% \vspace{-1.0\baselineskip}
\noindent
\begin{minipage}[t]{0.48\textwidth}\vspace{0pt}
\begin{itemize}
\item
\textbf{R5:} When a package is delivered to the selected locker, the customer receives a unique code (e.g., a 6-digit PIN) to open the locker.
\item
\textbf{R6:} Packages are held in the locker for three days.
\item
\textbf{R7:} Each locker location has defined opening and closing hours; customers must pick up packages within the 3-day window and the location's operating hours.
\end{itemize}
\end{minipage}
\hfill
\begin{minipage}[t]{0.48\textwidth}\vspace{0pt}
\textit{Column component type 'MxGraphWidget' not yet supported.}
\end{minipage}

% \vspace{-1.0\baselineskip}
\noindent
\begin{minipage}[t]{0.48\textwidth}\vspace{0pt}
\begin{itemize}
\item
\textbf{R8:} If a package is not picked up within three days, it is removed from the locker, the locker is released, and the customer is refunded.
\item
\textbf{R9:} Multiple lockers are available at every locker location; each can only be assigned to one customer/package at a time.
\item
\textbf{R10:} Once a package is collected, the locker is closed and locked; the provided access code is invalidated and cannot be reused.
\end{itemize}
\end{minipage}
\hfill
\begin{minipage}[t]{0.48\textwidth}\vspace{0pt}
\textit{Column component type 'MxGraphWidget' not yet supported.}
\end{minipage}

% \vspace{-1.0\baselineskip}
\noindent
\begin{minipage}[t]{0.48\textwidth}\vspace{0pt}
\begin{itemize}
\item
\textbf{R11:} Customers may return eligible items by selecting a nearby locker location. Based on package size and location, an available locker is assigned. A new, unique code is sent to the user to open the locker and place the return package.
\item
\textbf{R12:} The logistics team stores returned items in lockers for pickup; the logistics team uses a unique code to collect the returned package. The customer is notified once the return is processed, and the refund policy is applied per product eligibility.
\end{itemize}
\end{minipage}
\hfill
\begin{minipage}[t]{0.48\textwidth}\vspace{0pt}
\textit{Column component type 'MxGraphWidget' not yet supported.}
\end{minipage}

We've identified our requirements for the problem, and in the next lesson, we will define different use cases for the Amazon Locker system.

% End of section content